\section{Proof of local results}\label{section4}

A key reference for the study of inf\/initesimal automorphisms of parabolic geometries is \cite{Cap inf aut}. In that text, A.~\v{C}ap generalized to arbitrary parabolic geometries a bijective correspondence between conformal vector f\/ields and adjoint tractors (sections of the associated bundle to the canonical Cartan bundle, $\Gbdle \rightarrow M$, induced by the adjoint representation on $\g$) satisfying an identity involving the Cartan curvature, which was f\/irst discovered by A.R.~Gover in~\cite{Gover}. Moreover, the text of \v{C}ap relates this general bijective correspondence to the f\/irst splitting operator of a so-called curved BGG-sequence for the parabolic geometry, cf.\ Theorem~3.4 of~\cite{Cap inf aut}. The curvature identity of \cite{Cap inf aut} extends without dif\/f\/iculty to general inf\/initesimal automorphisms of Cartan geometries. This allows us to apply this fundamental identity to the Cartan geometries $(\Gbdle_0,\sigma^*\omega_{\leq})$, which we do in Section~\ref{section4.2} to establish a general ``dictionary'' between essentiality of an inf\/initesimal automorphism near a singularity, and the so-called holonomy associated to such a singularity (cf.\ Def\/inition~\ref{holonomy definition}; again, this should not be confused with the holonomy of the Cartan connection).

\subsection{Background results on inf\/initesimal automorphisms}\label{section4.1}

We recall some general notions, mainly following the development of \cite{Cap inf aut} (cf.\ also~1.5 of \cite{CSbook}), in the setting of a general Cartan geometry $(\Gbdle \rightarrow M,\omega)$ of type $(G,P)$ (for now not assumed to be of parabolic type). For any representation $\rho: P \rightarrow Gl(V)$, we have the associated vector bundle $V(M) := \Gbdle \times_{\rho} V$. The smooth sections of such a bundle are identif\/ied with $P$-equivariant, $V$-valued smooth functions on $\Gbdle$ in the standard manner, and we will simply treat them as such:
\[
\Gamma(V(M)) = \big\{f \in C^{\infty}(\Gbdle,V) \, \vert \, f(u.p) = \rho(p^{-1})(f(u))\big\} =: C^{\infty}(\Gbdle,V)^P.
\]
For the most part, the important associated bundles we are dealing with are \emph{tractor bundles}, which for our purposes simply means that the representation $(\rho,V)$ is the restriction to $P$ of a $G$-representation $\tilde{\rho}:G \rightarrow Gl(V)$. And the primary tractor bundle is the adjoint bundle induced by the restriction of the adjoint representation $\mathrm{Ad}: G \rightarrow Gl(\g)$ to $P$, which we will denote by $\Abdle = \Abdle(M)$ if there is no danger of confusion about which Lie algebra $\g$ is meant, and otherwise by $\g(M)$. Note that the Lie bracket $[\cdot \, ,\cdot]_{\g}$ of $\g$, by $\mathrm{Ad}(P)$-invariance, determines an algebraic bracket on f\/ibers of $\Abdle$ as well as on sections, which we denote with curly brackets $\{\cdot \, ,\cdot\}: \Abdle \times \Abdle \rightarrow \Abdle$. Also, note that there is a natural projection $\Pi: \Abdle \rightarrow TM$, induced by the projection $\g \rightarrow \g/\p$ and the isomorphism $TM \isom_{\omega} \Gbdle \times_{\overline{\mathrm{Ad}}(P)} \g/\p$.

The Cartan connection determines an identif\/ication of right-invariant vector f\/ields
\[
\VF(\Gbdle)^P = \{ \mathbf{X} \in \VF(\Gbdle) \, \vert \, \mathbf{X}(u.p) = (R_p)_* (\mathbf{X}(u))\},
\]
with sections of the adjoint bundle. Namely, to $\mathbf{X} \in \VF(\Gbdle)$ we associate a function $s_{\mathbf{X}} \in C^{\infty}(\Gbdle,\g)$ def\/ined by $s_{\mathbf{X}}(u) := \omega(\mathbf{X}(u))$; conversely, to a function $s \in C^{\infty}(\Gbdle,\g)$, associate $\mathbf{X}_s \in \VF(\Gbdle)$ def\/ined by $\mathbf{X}_s(u) = \omega_u^{-1}(s(u))$. The property~(\ref{parallelism}) of a~Cartan connection insures that both maps are well-def\/ined, they are inverse, and by property~(\ref{ad-inv}) of $\omega$ these maps restrict to an isomorphism $\VF(\Gbdle)^P \isom_{\omega} \Gamma(\Abdle)$. Similarly, the Cartan connection $\omega$ induces natural identif\/ications $\overline{\Omega^k(\Gbdle;\g)}^P \isom_{\omega} \Omega^k(M;\Abdle)$ of the horizontal, $\mathrm{Ad}(P)$-equivariant $\g$-valued $k$-forms on $\Gbdle$ with the $\Abdle$-valued $k$-forms on~$M$. For a tractor bundle~$V(M)$, the identif\/ication $\VF(\Gbdle)^P \isom_{\omega} \Gamma(\Abdle)$ yields two kinds of dif\/ferentiation of smooth sections with respect to adjoint tractors:

\begin{Definition} \label{differentiation} The \textit{invariant differentiation} or \textit{fundamental D-operator} of $V(M)$ is the map $D^{V} : \Gamma(V(M)) \rightarrow \Gamma(\Abdle^* \tens V(M))$ def\/ined, for any $s \in \Gamma(\Abdle)$ and any $v \in \Gamma(V(M))$, by:
\[
D_s^{V}v := \mathbf{X}_s(v).
\]
The \textit{tractor connection} of $V(M)$ is the map $\nabla^{V}: \Gamma(V(M)) \rightarrow \Gamma(\Abdle^* \tens V(M))$ def\/ined, for any $s \in \Gamma(\Abdle)$ and any $v \in \Gamma(V(M))$, by:
\begin{gather}
\nabla^{V}_s v := D^{V}_s v + (d\tilde{\rho} \circ s) \circ v. \label{tractor connection}
\end{gather}
 Recall that $\tilde{\rho}: G \rightarrow Gl(V)$ is the $G$-representation which~$\rho$ is a restriction of, given by the def\/inition of a tractor bundle. In fact, the quantity def\/ined by (\ref{tractor connection}) only depends on the equiva\-len\-ce class~$[s]$ of~$s$ under the quotient $\Abdle/\p(M) = \Gbdle \times_{\overline{\mathrm{Ad}}(P)} \g/\p \isom TM$, and we identify the tractor connection with a covariant derivative on~$V(M)$:
\[
\nabla^V: \ \Gamma(V(M)) \rightarrow \Gamma(T^*M \tens V(M)).
\]
\end{Definition}

The curvature tensor of a Cartan connection is the $\g$-valued two-form on $\Gbdle$ def\/ined, for any $u \in \Gbdle$ and $v, w \in T_u\Gbdle$, by the structure equation $\Omega^{\omega}(v,w) := d\omega(v,w) + [\omega(v),\omega(w)]$. The curvature tensor is horizontal and $P$-equivariant, and we may equivalently consider the curvature function $\kappa^{\omega} \in C^{\infty}(\Gbdle;\Lambda^2(\g/\p)^* \tens \g)^P$ induced by $\kappa^{\omega}(u)(X,Y) := \Omega^{\omega}(\omega_u^{-1}(X),\omega_u^{-1}(Y))$ for any $u \in \Gbdle$ and $X,Y \in \g$. The following identity was proven for parabolic geometries in~\cite{Cap inf aut} and the proof carries over without substantive changes to Cartan geometries of general type (cf.\ also Lemma~1.5.12 in~\cite{CSbook}):

\begin{Lemma} \label{infinitesimal deformation} Let $\mathbf{X} \simeq_{\omega} s_{\mathbf{X}}$ for $\mathbf{X} \in \VF(\Gbdle)^P$ and $s_{\mathbf{X}} \in \Gamma(\Abdle)$. Then for $\mathcal{L}_{\mathbf{X}} \omega \in \overline{\Omega^1(\Gbdle;\g)}^P$ we have, under the identification $\overline{\Omega^1(\Gbdle;\g)}^P \isom_{\omega} \Omega^1(M;\Abdle)$:
\begin{gather}
\mathcal{L}_{\mathbf{X}} \omega \simeq_{\omega} \nabla^{\Abdle} s_{\mathbf{X}} + \Pi(s_{\mathbf{X}}) \: \k^{\omega}. \label{infinitesimal deformation identity}
\end{gather}
\end{Lemma}

\subsection{Holonomy and essentiality of inf\/initesimal automorphisms}\label{section4.2}

Let us return now to the setting of a (regular, normal) parabolic geometry $(\Gbdle \rightarrow M,\omega)$ of type $(G,P)$ and the corresponding regular inf\/initesimal f\/lag structure $\mathcal{M} = (M,\{T^iM\},\Gbdle_0)$ of type $(\g,P)$. As was noted for automorphisms of $(\Gbdle,\omega^{nc})$ in Remark \ref{underlying automorphisms}, we may determine an inf\/initesimal automorphism $\mathbf{X} \in \mathrm{inf}(\Gbdle,\omega)$ by conditions on the underlying vector f\/ield $X \in \VF(M)$, which just amount to imposing the same conditions for the locally def\/ined dif\/feomorphisms given by f\/lowing along $X$, i.e.\ we must have $[X,\Gamma(T^iM)] \subseteq \Gamma(T^iM)$ for all $-k \leq i \leq -1$, and the condition that the local f\/lows of $X$ determine local bundle maps of $\mathcal{F}(\mathrm{gr}(TM))$ which preserve~$\Gbdle_0$ as a~subbundle. In the dif\/ferent examples of parabolic geometries, this translates into more geometric language. For example, in the conformal case, the former condition is trivial, while the latter condition amounts to requiring the conformal Killing equation, $\Lie_X g = \lambda g$ for any $g \in c$ and some $\lambda = \lambda(X) \in C^{\infty}(M)$. In the case of CR structures, the conditions are that $[X,\Gamma(\mathcal{H})] \subseteq \Gamma(\mathcal{H})$ for $\mathcal{H} \subset TM$ the codimension one contact distribution def\/ining the CR structure, and that $\Lie_X J = 0$ for $J$ the almost complex structure on $\mathcal{H}$. We write $X \in \mathrm{inf}(\mathcal{M})$ and consider the lift $\mathbf{X} \in \mathrm{inf}(\Gbdle,\omega^{nc})$ to be implicitly included.

Using Lemma~\ref{infinitesimal deformation} and Lemma~\ref{inessential automorphisms TFAE}, we obtain a bijection between inf\/initesimal automorphisms $\mathbf{X} \in \mathrm{inf}(\Gbdle,\omega)$ and adjoint tractors $s_{\mathbf{X}} \in \Gamma(\Abdle)$ satisfying $\nabla^{\Abdle}s_{\mathbf{X}} + \Pi(s_{\mathbf{X}}) \: \k^{\omega} = 0$. In the present setting, this gives us a bijection between $X \in \mathrm{inf}(\mathcal{M})$ and such $s_{\mathbf{X}}$, and moreover it is easy to verify that $\Pi(s_{\mathbf{X}}) = X \in \VF(M)$. Now, denote by $\mathbf{X}_0 \in \VF(\Gbdle_0)^{G_0}$ the invariant vector f\/ield induced, via projection by $\pi_+$, by $\mathbf{X} \in \VF(\Gbdle)^P$. For any Weyl structure $\sigma: \Gbdle_0 \rightarrow \Gbdle$, the induced Cartan connection $\sigma^*\omega_{\leq}$ gives us an isomorphism denoted $\VF(\Gbdle_0)^{G_0} \isom_{\sigma} \Gamma(\Abdle^{\sigma})$, where we denote with $\Abdle^{\sigma} = \p^*(M)$ the adjoint bundle of the Cartan geometry $(\Gbdle_0 \rightarrow M,\sigma^*\omega_{\leq})$. Let us write $\mathbf{X}_0 \simeq_{\sigma} s_{\mathbf{X}_0}$ for the adjoint tractor corresponding to $\mathbf{X}_0 \in \VF(\Gbdle_0)^{G_0}$. If we further denote by $\nabla^{\sigma}: \Gamma(\Abdle^{\sigma}) \rightarrow \Gamma(T^*M\tens\Abdle^{\sigma})$ the corresponding adjoint tractor connection, then Lemma \ref{infinitesimal deformation} tells us that $X \in \mathrm{inf}(\mathcal{M})$ is an inf\/initesimal automorphism of some Weyl structure $\sigma: \Gbdle_0 \rightarrow \Gbdle$ if and only if
\[
\nabla^{\sigma} s_{\mathbf{X}_0} + X \: \k^{\sigma^*\omega_{\leq}} = 0.
\]

At present, we are only interested in local properties of inf\/initesimal automorphisms, viz the question if some neighborhood of a given point can be found on which $X$ is inessential. The following proposition shows that the points $x \in M$ for which the answer could be ``no'', must all be singularities of the vector f\/ield, i.e.\ $X(x)=0$ (where, incidentally, the above identity simplif\/ies to $(\nabla^{\sigma} s_{\mathbf{X}_0})(x) = 0$).
